\documentclass[11pt]{article}

\usepackage[a4paper,margin=2.5cm]{geometry}
\usepackage[T1]{fontenc}
\usepackage[utf8]{inputenc}
\usepackage{lmodern}
\usepackage{xcolor}
\usepackage{amsmath,amssymb}
\usepackage[hidelinks]{hyperref}

\setlength{\parindent}{0pt}
\setlength{\parskip}{0.7em}

\definecolor{reviewgray}{RGB}{70,70,70}
\definecolor{responseblue}{RGB}{0,72,125}

\newenvironment{reviewercomment}{\begin{quote}\color{reviewgray}\itshape}{\end{quote}}
\newenvironment{authorresponse}{\par\noindent\textbf{\color{responseblue}Response.}\ }{\par}
\newcommand{\change}[1]{\par\noindent\textit{Changes in the revised manuscript.} #1\par}

\begin{document}

\begin{center}
{\Large\bfseries Response to Reviewer 1}\\[0.5em]
{\itshape Manuscript: ``An attractor account of numerosity perception''}\\[0.25em]
{\itshape Manuscript ID: IPE-26-0055}
\end{center}

Dear Reviewer,

Thank you for the careful, technically detailed assessment. Your comments led us to rebuild the principal likelihood and model comparison, derive the response-time prediction from the dynamics, test a landscape shared across exposure durations, and narrow several theoretical claims. In short, what began as a small project exploring a new idea has become a substantially more rigorous computational cognition paper, and I am grateful for the care you took in this review.

We address every major and minor point below.

\subsection*{Major Comment 1: The model comparison used the wrong likelihood}

\begin{reviewercomment}
The original AIC/BIC comparison was computed from Gaussian residuals on three summary curves. The reviewer requested a trial-level response likelihood for the discrete integer reports, with response time modelled separately.
\end{reviewercomment}

\begin{authorresponse}
We agree. The original comparison discarded important information in the discrete response distributions and did not provide an appropriate basis for the principal model comparison.

We therefore rebuilt the analysis around a categorical trial-level likelihood. For each presented numerosity, the probability of every integer response is obtained by marginalising the model readout over the sensory distribution,
\[
p(r\mid n)=\int p\!\left(r\mid x^*(y_0)\right)p(y_0\mid n)\,dy_0.
\]
The integral is evaluated deterministically, and the likelihood is computed over all 52,995 individual integer reports. Typed responses outside the integer-well support are retained through the same uniform lapse channel in every model, and the lapse is included in all AIC and BIC calculations.

Response times are no longer combined with response summaries in a single loss. They are modelled separately using a trial-level log-normal likelihood, with the landscape fixed at the values obtained from the categorical-response fit. The response-time analysis is therefore a constrained cross-channel test of the response-fitted mechanism.
\end{authorresponse}

\change{The revised Methods and Supplementary Material~1 specify the categorical likelihood, deterministic quadrature, lapse convention, well and response supports, and separate response-time channel. The model comparison is reported in Figure~2a and Supplementary Table~S2.}

\subsection*{Major Comment 2: The reported AIC differences were too small}

\begin{reviewercomment}
The original summary-level difference between the Hopfield--Poisson model and its no-clean-up control was only $\Delta\mathrm{AIC}=1.6$, which did not support strong claims about the contribution of the dynamics.
\end{reviewercomment}

\begin{authorresponse}
We agree that the original summary-level differences did not support strong conclusions about model preference. This was one of the principal motivations for replacing the original analysis.

Under the revised trial-level likelihood, relative to the power-law attractor, the matched static power-law readout trails by $\Delta\mathrm{AIC}=722$, the constant-width attractor by $\Delta\mathrm{AIC}=3488$, the Weber/ANS model by $\Delta\mathrm{AIC}=838$, the matched bounded-optimal model by $\Delta\mathrm{AIC}=11{,}453$, and the purely driven accumulator by $\Delta\mathrm{AIC}=7{,}954$. BIC gives the same model ordering.

Importantly, the static control has one fewer parameter because it does not require the shared stopping budget $G$. The attractor advantage therefore remains after penalizing this additional flexibility.
\end{authorresponse}

\change{The manuscript now bases its empirical case on the trial-level comparison in Figure~2a and Supplementary Table~S2 rather than on the original summary-curve AIC values.}

\subsection*{Major Comment 3: The Hopfield framing was largely decorative}

\begin{reviewercomment}
The reviewer questioned what the one-dimensional attractor dynamics contributed beyond a static mixture or noisy MAP readout, and noted that the model did not address the high-dimensional storage-capacity properties usually associated with modern Hopfield networks.
\end{reviewercomment}

\begin{authorresponse}
We agree that the original analysis did not adequately isolate what the attractor dynamics contributed beyond the static response family.

The revised manuscript contains a matched static control that retains the same sensory channel, integer-centred Gaussian components, width law, and categorical readout but removes gradient-descent clean-up. Despite requiring one fewer parameter, this static model fits substantially worse than the attractor model. We also include constant-width versions of the static and attractor models, which separately test whether the advantage comes from the dynamics or simply from magnitude-dependent broadening.

The dynamics make a second distinctive prediction through convergence time. The selected response-time model maps the expected number of bounded clean-up steps onto response time; an entropy-only model performs substantially worse. Thus, the chronometric analysis depends directly on the operation that differentiates the attractor from its static twin.

We have also narrowed the theoretical use of the Hopfield framework. We do not claim the high-dimensional storage capacity associated with modern Hopfield networks. The relevant claim is a one-dimensional associative-retrieval operation in which a noisy scalar state is iteratively transformed by a log-sum-exp energy landscape generated by stored integer-centred components, in a modern Hopfield view, but it's a particular case where the whole attractor network is only one dimensional. The revised manuscript therefore describes the model as Hopfield-inspired and focuses on pattern completion rather than capacity; ``Hopfield'' has also been removed from the title, forgoing a bit of flair in favour of greater precision.
\end{authorresponse}

\change{The matched static and constant-width controls are shown in Figure~2a and Supplementary Table~S2; the convergence-time analysis is shown in Figures~1c and~2b. The revised Methods and Discussion delimit the Hopfield analogy explicitly.}

\subsection*{Major Comment 4: The accumulator interpretation was not derived}

\begin{reviewercomment}
The reviewer argued that the original Poisson wells were a parametric choice rather than a genuine implementation of a classical accumulator, and that the manuscript used the accumulator label more strongly than the formal model justified.
\end{reviewercomment}

\begin{authorresponse}
We agree. We have removed the original Poisson-well parameterisation and no longer describe the main model as an accumulator model.

The revised landscape is generated by Gaussian memory components centred on the integers, with widths
\[
\sigma_n=\sigma_c n^p.
\]
The power-law broadening is motivated by the decline in number-word frequency, rather than presented as a derivation from accumulation statistics. The model now begins with an already formed scalar numerosity estimate $y_0$, and the attractor implements subsequent clean-up and numerical readout. We also compared this model with the earlier Poisson-shaped landscape; it fitted substantially worse, and we removed it because we lacked a theoretical basis for retaining Poisson-shaped wells.

On the other hand, we did model a purely accumulator version of the attractor to correctly test the original accumulator intuition explicitly. In a purely driven model, registered items generate a running count $K(t)$ that drives the attractor during exposure. Alas, this model performed substantially worse than the completed-estimate attractor, so we argue in the manuscript against a unified process of numerical perception (in the sense of accumulation and readout being two processes and not one). We argue that this two-step process makes sense considering that a partial count is itself a valid intermediate cardinality and is not diagnostic of the final total, so the attractor is repeatedly drawn towards intermediate integer memories. We now interpret the attractor as a mechanism operating on the output of accumulation rather than as the accumulation process itself.

\end{authorresponse}

\change{The Poisson formulation and accumulator claims have been removed. The purely driven running-count model is now a quantitative control described in Methods, Results, Discussion, and Supplementary Figure~S1.}

\subsection*{Major Comment 5: Response time was not derived from the dynamics}

\begin{reviewercomment}
The original response-time model used posterior entropy. The reviewer suggested deriving response time from convergence time or first-passage time on the attractor landscape.
\end{reviewercomment}

\begin{authorresponse}
We agree with this assessment and have replaced the original entropy-based account. This was also our initial intuition; however, we found that for large numerosities, the attractor mechanism fails to converge within a plausible number of steps. Consequently, we had to introduce a "give-up" mechanism to bound this number of steps, otherwise we got extremely large response times. We initially considered the introduction of this unexplained mechanism to be a "Deus ex machina" solution, leading us to favour the classical entropy readout. Nonetheless, following your comment, we have revisited our decision and reintroduced it. When compared to the entropy readout, this version of the model is far superior; we therefore favour it and stand by the "give-up" mechanism, which constitutes a prediction of the model that we will further analyse in future work.
\end{authorresponse}

\change{The integration rule, convergence and oscillation criteria, step clamp, stopping budget, empirical response-time origin, log-normal likelihood, and fitted response-time parameters are now documented in Methods and Supplementary Material~1.}

\subsection*{Major Comment 6: The three-regime and subitizing claims were overstated}

\begin{reviewercomment}
The reviewer noted that the density regime was not modelled and argued that fitting small-number accuracy did not amount to a mechanistic account of all classical subitizing phenomena, because the model contains no spatial, object-based, or attentional mechanism.
\end{reviewercomment}

\begin{authorresponse}
We agree that the original manuscript conflated three distinct claims: fitting small-number behaviour, providing a complete mechanism of subitizing, and predicting a transition to density perception.

We have removed the density regime entirely, including the proposed density transition and the comparison with Anobile et al. The revised paper concerns the change from more categorical to increasingly approximate numerical reports within the fitted range.

We have also narrowed the subitizing claim. Under the revised trial-level likelihood, the model reproduces near-ceiling exact-response rates for the smallest numerosities, but we agree that this does not constitute a model of all classical subitizing phenomena. The attractor has no representation of object locations, configurations, or attentional load. We therefore describe it as a model of numerical readout downstream of modality-specific sensory and attentional processing, and state that these perceptual contributions lie outside its present scope.

We retain a narrower unification claim: the same numerical readout produces more categorical reports where integer components are well separated and Weber-like reports where they overlap. This does not require every perceptual signature of subitizing to arise from the landscape, nor does it rule out object-individuation processes for the smallest sets.
\end{authorresponse}

\change{All density-regime analyses and claims have been removed. The Abstract, Introduction, Results, and Discussion now restrict the proposal to a shared numerical readout, and the Discussion identifies spatial, continuous-magnitude, and attentional effects as outside the present model.}

\subsection*{Major Comment 7: Insufficient engagement with competing literature}

\subsubsection*{Cheyette, Wu, and Piantadosi (2024)}

\begin{reviewercomment}
The reviewer asked us to engage directly with the visuospatial-capacity account, which attributes number psychophysics to bandwidth-limited visual memory and reports within-participant transfer between non-numerical visual memory and numerosity estimation.
\end{reviewercomment}

\begin{authorresponse}
We agree that the original treatment was insufficient. The revised manuscript now presents the attractor and visuospatial-capacity accounts as competing explanations of where the same psychophysical regularities arise.

The visuospatial account places the principal limitation in sensory encoding. We now discuss the within-participant transfer of information-capacity estimates from change localization to number estimation reported by Cheyette et al. (2024), which provides direct evidence for a shared visuospatial constraint.

The attractor account instead predicts that a shared numerical landscape should generalize across sensory modalities and may be selectively modified by numerical experience. We formulate differential tests involving cross-modal parameter sharing and numerical versus visuospatial training. We also discuss evidence that nonsymbolic acuity can remain robust with little formal schooling, which constrains a strongly schooling-dependent interpretation of the landscape.

A hierarchical individual-differences analysis provides an additional test within the present account. It uses all five exposure durations and non-overlapping numerosity ranges: $n\leq8$ informs participant landscape resolution and $n\geq9$ informs Weber noise. Participant effects on both outcomes and on lapse are jointly modelled, with the lapse effect shared across durations and response channels. Greater landscape resolution predicts lower Weber noise conditional on lapse ($\hat\rho_{nw\cdot\ell}=-.627$, $N=360$, $\chi^2(1)=105.0$, $p=1.2\times10^{-24}$). We describe this as a cross-participant association, not causal evidence.
\end{authorresponse}

\change{The competing explanatory loci and their differential predictions are now discussed in the Introduction and Discussion; the all-duration individual-differences analysis is reported in Figure~2c and Supplementary Material~4.}

\subsubsection*{Recurrent visual-network literature}

\begin{reviewercomment}
The reviewer asked us to engage with the recurrent on-centre/off-surround models developed by Sengupta and colleagues, including their range effects, adaptation, and switch-cost predictions.
\end{reviewercomment}

\begin{authorresponse}
We have expanded the Introduction to acknowledge the recurrent visual-network literature, including Sengupta et al. (2014), Sengupta et al. (2017), and Verma and Sengupta (2023), and we no longer claim novelty for a three-range or density boundary.

We distinguish these accounts from the present proposal by explanatory level. The recurrent visual models derive numerical sensitivity, range effects, adaptation, and switch costs from the dynamics of a sensory network, whereas the present model begins with scalar numerical evidence and asks how that evidence is stabilized by an explicit memory landscape. The latter phenomena are important tests not addressed by the present dataset.
\end{authorresponse}

\change{The recurrent-network literature and the distinction between sensory network dynamics and central numerical readout have been added to the Introduction; the Discussion identifies a recurrent visual front end as a natural extension.}

\subsection*{Major Comment 8: Independent fits under-tested model stability}

\begin{reviewercomment}
The reviewer suggested a constrained model in which the representational landscape is shared across exposure durations and only sensory parameters vary.
\end{reviewercomment}

\begin{authorresponse}
We agree. The primary analysis is now a constrained joint model with a single power-law landscape shared across all five exposure durations. Specifically, $\sigma_c$, $p$, and the global stopping budget $G$ are shared, while the sensory-noise parameters and lapse vary with exposure duration.

This is now the headline model throughout the manuscript. It provides the best trial-level account despite being more constrained than the fully duration-specific Weber/ANS model. The shared-landscape attractor beats the Weber/ANS model by $\Delta\mathrm{AIC}=838$ and the matched bounded-optimal model by $\Delta\mathrm{AIC}=11{,}453$.

We therefore no longer infer that the cognitive system uses a different numerical landscape at each duration. Exposure duration is represented through the sensory channel, while the numerical landscape is treated as a stable representational property.
\end{authorresponse}

\change{The sharing structure and parameter counts are specified in Methods, Supplementary Material~1, and Supplementary Table~S2.}

\subsection*{Major Comment 9: The dyscalculia extension was too weak}

\begin{reviewercomment}
The reviewer argued that the original group-level analysis and its interpretation were too strong for the small dataset and limited numerosity range, and questioned the exclusion of $n=8$.
\end{reviewercomment}

\begin{authorresponse}
We agree that the original group-level fit could not support the causal interpretation we gave it, and we have removed it.
\end{authorresponse}

\change{Supplementary Material~5 and all corresponding main-text references to the exploratory dyscalculia analysis have been removed.}

\subsection*{Minor Comment: Benford framing}

\begin{reviewercomment}
Dehaene and Mehler (1992) report a power-law decline in number-word frequency, not strict Benford behaviour.
\end{reviewercomment}

\begin{authorresponse}
We agree. All references to Benford's law have been removed. The revised model uses a simple power-law width schedule motivated by the cross-linguistic power-law decline in number-word frequency. We do not claim that the empirical frequencies follow strict Benford behaviour.
\end{authorresponse}

\change{The motivation and width law have been rewritten in the Introduction, Methods, and Supplementary Material~2.}

\subsection*{Minor Comment: The 40-ms condition}

\begin{reviewercomment}
The reviewer requested a more candid account of what fails at 40 ms rather than describing the condition only as noisier.
\end{reviewercomment}

\begin{authorresponse}
We agree. The revised manuscript now identifies the 40-ms condition as one of the two systematic residuals. It accumulates more deviations from the fitted mean-estimate, exact-accuracy, and response-time curves than the longer exposures; the other residual is a model-predicted shoulder that is less evident in participants' accuracy and response times. We report these deviations directly rather than treating the shortest exposure merely as noisy data.
\end{authorresponse}

\change{Both systematic residuals are now described in the Results alongside Figure~3.}

\subsection*{Minor Comment: Comparison with Anobile et al.}

\begin{reviewercomment}
The reviewer regarded the comparison between exposure-dependent model thresholds and patch-size-dependent density thresholds as numerological.
\end{reviewercomment}

\begin{authorresponse}
We agree. The proposed density transition and the comparison with Anobile et al. have been removed entirely. The revised manuscript does not model density perception.
\end{authorresponse}

\change{The density analysis, figure, and associated claims have been removed from the main text and Supplementary Material.}

\subsection*{Minor Comment: Compression and logarithmic geometry}

\begin{reviewercomment}
The reviewer noted that the original model did not adjudicate between increasing well width on a linear axis and fixed-width wells on a logarithmically compressed axis.
\end{reviewercomment}

\begin{authorresponse}
We agree. The Poisson accumulator parameterization and the claim that it established the origin of compression have been removed. The revised paper tests a power-law width schedule on a linear axis but does not treat that fit as uniquely identifying number-line geometry. An appropriately compressed coordinate system could in principle provide a related description; we have not performed an exhaustive comparison of alternative coordinate geometries.
\end{authorresponse}

\change{The original compression claim has been removed, and the Discussion now states this geometric non-identifiability explicitly.}

We thank the reviewer again. The revised paper is substantially different from the original submission: its principal evidence now comes from a trial-level, shared-landscape analysis; the dynamics are isolated by matched controls and linked directly to convergence time; and the theoretical claims are restricted to the level supported by the data.

Sincerely,

Eudald Correig-Fraga

\end{document}
